\section{Signed phase geometry: disks and exact circles}\label{sec:phases}

\subsection{The full quadratic form}

Suppress the subscript on \(K_a\), and write arbitrary real quadratures
\[
 x_n=d_n\cos\phi_n,\qquad y_n=d_n\sin\phi_n,\qquad
 q(w)=\sum_{n\le N}[x_n\cos(\lambda_nw)-y_n\sin(\lambda_nw)]+K(w).
\]
The actual magnitudes obey \(x_n^2+y_n^2=d_n^2\); a disk relaxation
replaces equality by \(\le\). The two event equations read
\begin{equation}\label{eq:eventquadratures}
 q(0)=\sum_nx_n+K(0)=0,\qquad
 q'(0)=-\sum_n\lambda_n y_n+K'(0)=0.
\end{equation}
The corresponding approximate class allows errors \(O(a^{-3/4})\).
Every exact solution belongs to that approximate class.

At \(\ii\beta\), write \(S=q^\circ\) as \(S(\ii\beta)=X-\ii Y\), where
\begin{align}
 C_n&=\cosh(\beta\lambda_n)-1,&
 B_n&=\sinh(\beta\lambda_n)-\beta\lambda_n,\\
 R_n&=\lambda_n\sinh(\beta\lambda_n),&
 V_n&=\lambda_n(\cosh(\beta\lambda_n)-1),\\
 X&=C\cdot x+\Rea K^\circ(\ii\beta),&
 Y&=B\cdot y-\Ima K^\circ(\ii\beta).
\end{align}
With primes on \(X,Y\) denoting \(\beta\)-derivatives,
\begin{equation}\label{eq:XY}
 \cD_a(S)=\mathfrak d_a(X^2+Y^2)-\beta(XX'+YY').
\end{equation}
Thus, apart from the explicitly retained correction terms, its matrices
are
\[
 Q_x=\mathfrak d_a CC^t-\frac\beta2(CR^t+RC^t),\qquad
 Q_y=\mathfrak d_a BB^t-\frac\beta2(BV^t+VB^t).
\]
These formulas keep all cross terms. They will allow counterexamples
for the full form, independently of any numerical optimization method.
Write \(\cD_2\) for the same functional with coefficient \(2\) in
place of \(\mathfrak d_a\).

\subsection{A two-mode disk construction}

Let
\[
 r=\frac14,\qquad h=\frac{\log2}{2L},\qquad t=\frac{d_2}{2}.
\]
Choose
\[
 x_1=t-K(0),\quad x_2=-t,\quad
 y_1=4K'(0),\quad y_2=0,
\]
and set every other quadrature to zero. Since
\(2^{-7/4}\le t\le2^{-3/2}\) and \(K,K'=O(\chi)\), these points
lie strictly inside their coefficient disks for large \(a\).
Both equations \eqref{eq:eventquadratures} are exact. The resulting
field is
\begin{equation}\label{eq:diskfield}
 q_*(w)=t\{\cos(rw)-\cos((r-h)w)\}+J(w),
 \quad J=K-K(0)\cos(rw)-4K'(0)\sin(rw).
\end{equation}
Its correction \(J\) and fixed derivatives are \(O(\chi)\).

Define the positive number
\begin{equation}\label{eq:Cbeta}
 C_\beta=\frac{(\beta\log2)^2}{4}\sinh^2(\beta/4)
             \left[\frac\beta4\coth(\beta/4)-1\right]>0.
\end{equation}
Positivity follows from \(x\cosh x-\sinh x>0\) for \(x>0\).
For \(U=\cosh(\beta r)-\cosh(\beta(r-h))\), expansion in \(h\) gives
\[
 U=\beta h\sinh(\beta r)+O_\beta(h^2),\qquad
 U(\beta U'-2U)=C_\beta/L^2+O_\beta(L^{-3}).
\]
The uncorrected field is real at \(\ii\beta\), so
\(\cD_2=-t^2U(\beta U'-2U)\). Expanding the entire correction and
carrier contribution yields
\begin{equation}\label{eq:disknegative}
 \cD_a(q_*)=-\frac{t^2C_\beta}{L^2}
   +O_\beta(L^{-3}+\chi/L+\chi^2+a^{-2}L^{-4})<0
\end{equation}
eventually. For example it is at most
\(-2^{-9/2}C_\beta/L^2\) after a finite onset.
This disproves nonnegativity on coefficient disks subject to the two
event equations. Exact magnitudes and the common arithmetic phase law
are absent from the witness.

\subsection{An elementary three-circle lemma}

The next construction restores every magnitude. The following lemma
will also be the local building block for multiplicative phases.

\begin{lemma}[A signed block on three circles]\label{lem:triple}
Let \(b_1>b_2>b_3>0\), \(b_1<b_2+b_3\), and
\(r_1>r_2>r_3\ge0\), with \(r_2>0\). Put \(\rho=r_3/r_2\).
There is a unique \(v\in(0,b_3)\) such that
\[
 A_2=\sqrt{b_2^2-\rho^2v^2},\qquad
 A_3=\sqrt{b_3^2-v^2},\qquad A_2+A_3=b_1.
\]
The quadratures
\begin{equation}\label{eq:triple}
 (x_1,y_1)=(b_1,0),\quad
 (x_2,y_2)=(-A_2,\rho v),\quad
 (x_3,y_3)=(-A_3,-v)
\end{equation}
have exact magnitudes, \(\sum x_i=\sum r_i y_i=0\).
Their field \(p(w)=\sum_i(x_i\cos r_iw-y_i\sin r_iw)\) satisfies
\(p(0)=p'(0)=0\) and \(\cD_2(p)<0\) for every \(\beta>0\).
Any finite sum of these blocks has the same strict sign.
\end{lemma}
\begin{proof}
The sum of the two square roots is continuous and strictly decreasing
on \([0,b_3]\). Its initial value is \(b_2+b_3>b_1\); its final value
is at most \(b_2<b_1\). This proves the unique interior solution.
The circle identities and event equations follow directly.

Writing \(p(\ii\beta)=X-\ii Y\) gives
\begin{align}
 X={}&A_2[\cosh(\beta r_1)-\cosh(\beta r_2)]
      +A_3[\cosh(\beta r_1)-\cosh(\beta r_3)]>0,\\
 Y={}&vr_3\left[\frac{\sinh(\beta r_2)}{r_2}
                       -\frac{\sinh(\beta r_3)}{r_3}\right]\ge0.
 \label{eq:tripleXY}
\end{align}
Use the continuous interpretation at \(r_3=0\). The power series for
\(X\) has positive coefficients in even powers starting with
\(\beta^2\). Consequently \(\beta X'-2X>0\), since the quartic
coefficient is positive. The series for \(Y\) has nonnegative
coefficients in odd powers starting with \(\beta^3\), so
\(\beta Y'-2Y\ge0\).
Equation~\eqref{eq:XY} becomes
\[
 \cD_2(p)=-X(\beta X'-2X)-Y(\beta Y'-2Y)<0.
\]
For sums of blocks the four quantities retain the same signs.
Thus the cross terms between blocks also have the required sign.
\end{proof}

\subsection{All coefficients on their circles}

Partition \(1,\ldots,N\) into consecutive triples, leaving at most two
terminal terms. For a triple \(j,j+1,j+2\), set
\(b_i=d_{j+i-1}\), \(r_i=\lambda_{j+i-1}\).
The logarithmic slope from Lemma~\ref{lem:moments} gives
\[
 b_2+b_3\ge(2^{-3/4}+3^{-3/4})b_1>b_1.
\]
Lemma~\ref{lem:triple} applies, including a possible last frequency
zero. Let \(p\) be the sum of all complete blocks.
Frequency differences in a block are \(O(1/(jL))\).
The event cancellations, applied to the cosine and sine differences
in \eqref{eq:tripleXY} and their analytic versions, give
\[
 \norm{p_j^{(k)}}_{|w|\le R}\le C_R\frac{d_j}{jL}
                 \le \frac{C_R}{Lj^{3/2}},\quad 0\le k\le2.
\]
The possible small \(r_3\) is harmless: use the analytic function
\(\sin(rw)/r\) in the odd term. Summing the convergent \(j^{-3/2}\)
bound yields
\begin{equation}\label{eq:circlefieldbound}
 \norm{p^{(k)}}_{|w|\le R}=O_R(L^{-1}),\qquad 0\le k\le2.
\end{equation}

The head block gives a quantitative negative margin. With
\(\alpha=1-1/\sqrt3>0\), its \(A_2=1-A_3\ge1-d_3\ge\alpha\).
The two first frequencies are \(r=1/4\) and \(r-h\).
All signs in Lemma~\ref{lem:triple} align, so \eqref{eq:Cbeta} implies
\begin{equation}\label{eq:circleleading}
 \cD_2(p)\le-\alpha^2 C_\beta/L^2+O_\beta(L^{-3}).
\end{equation}

Give each leftover terminal term phase zero and denote their sum by
\(b_{\rm tail}(w)\). Since, for \(\ell=0,1\),
\[
 d_{N-\ell}
 =\exp[-(1/4+c/16)L+O(\ee^{-L/2})]=O(\chi),
\]
this tail and its fixed derivatives are \(O_R(\chi)\), and
\(b_{\rm tail}'(0)=0\).
It remains to restore the two event equations in the presence of the
complete \(K\), without leaving a circle.

Vary just the second and third phases of the head block. Their event
map is
\[
 F(\theta_2,\theta_3)=
 \bigl(1+d_2\cos\theta_2+d_3\cos\theta_3,\
       \lambda_2d_2\sin\theta_2+\lambda_3d_3\sin\theta_3\bigr).
\]
At the block solution its Jacobian determinant is
\[
 v\left[\frac{\lambda_3^2}{\lambda_2}A_3+\lambda_2A_2\right]>0.
\]
As \(L\to\infty\), the two frequencies tend to \(1/4\), and
\((1,d_2,d_3)\) tends uniformly to \((1,2^{-s},3^{-s})\) for
\(s\in[5/8,3/4]\). These are a compact family of nondegenerate
triangles, so the inverse Jacobian and the second derivatives of the
map are uniformly bounded. The inverse-function theorem therefore
solves
\[
 F=\bigl(-K(0)-b_{\rm tail}(0),\,K'(0)\bigr)
\]
with phase changes \(O(\chi)\). The complete repaired field
\(q_*=p+E\) satisfies the event equations exactly and
\(\norm{E^{(k)}}_R=O_R(\chi)\) for \(k\le2\).

\begin{theorem}[Exact-circle obstruction]\label{thm:circles}
For each fixed \(\beta>0\), every sufficiently large center in
\eqref{eq:band} admits quadratures with
\(x_n^2+y_n^2=d_n^2\) for every \(n\le N\), the actual correction \(K\),
and \(q_*(0)=q_*'(0)=0\), such that
\begin{equation}\label{eq:circlenegative}
 \cD_a(q_*)\le-\kappa_\beta^{\rm circ}/L^2,\qquad
 \kappa_\beta^{\rm circ}=\tfrac12(1-1/\sqrt3)^2C_\beta>0.
\end{equation}
\end{theorem}
\begin{proof}
Expand \(\cD_a(p+E)\) using \eqref{eq:circlefieldbound}.
The error beyond \eqref{eq:circleleading} is
\[
 O_\beta(\chi/L+\chi^2+a^{-2}L^{-4}).
\]
It is \(o(L^{-2})\), proving the result.
\end{proof}

This exact-circle point is still not multiplicative. The first member
of each block has phase zero, so \(u_1=u_4=1\). The adjusted second
phase remains uniformly in the interior of \((\pi/2,\pi)\). Therefore
\[
 |u_4u_1-u_2^2|=2|\sin\theta_2|\ge\delta_0>0.
\]
Nor is the actual leading carrier phase imposed in this construction.
The next section restores both multiplicativity and that leading phase.
